\section{Vectorized Kernel Details}
\label{app:kernels}

This appendix details the architecture-specific cores summarized in
\Cref{sec:impl}: the instruction selection, register layout, and
remainder handling of the \texttt{amd64} and \texttt{arm64} kernels.

\paragraph{\texttt{amd64}.}
The \texttt{amd64} implementation provides AVX-512 and AVX2 cores and selects
between them once at startup from the \texttt{AVX512VL} CPUID feature bit. The
single-duplex permutation has two variants: a general-purpose-register core
adapted from the standard library's \Keccak{} permutation (using the
lane-complement trick so that $\chi$ requires no separate negation), and an
AVX-512 variant; the root duplex uses whichever matches the selected core. The
eight-way AVX-512 core packs the lane-major state into 25 ZMM registers, one per
lane with all eight instances in the register's eight 64-bit slots, and performs
the permutation with no cross-lane shuffles: rotations use \texttt{VPROLQ}, and
\texttt{VPTERNLOGQ} fuses the $\theta$ $D$-term formation and the $\chi$
\mbox{and-not} into single three-input logic instructions. The eight-way AVX2
core, lacking 512-bit registers, runs the eight instances as two groups of four
(four 64-bit lanes per YMM register), rotating by shift-or and computing $\chi$
with \texttt{VPANDN}. The fused chunk kernels combine absorption with the
permutation. The AVX-512 steady-state kernel reads the eight chunks with
contiguous loads, transposes each rate block into the lane-major layout in
registers (\texttt{VPUNPCKLQDQ}/\texttt{VPUNPCKHQDQ} and \texttt{VSHUFI64X2}),
absorbs on the lane-major state, and transposes each output block back for a
contiguous store, leaving only the $7$-byte partial rate lane to a masked gather;
these sequential accesses are prefetcher-friendly, whereas per-lane gathers and
scatters at the chunk stride $\beta$
(\texttt{VPGATHERQQ}/\texttt{VPSCATTERQQ}) would stall on cold memory and cap
the long-message rate. The register-resident kernel that drains a
two-to-seven-chunk remainder instead uses that per-lane gather/scatter form
under a mask, reading the active chunks in place: the transpose's shuffle work
is constant regardless of the chunk count and would be spent partly on inactive
lanes, whereas the gathers scale with the active element count. The AVX2 path
drains the same remainders with a variant of its grouped-by-four kernels that
aims each inactive instance's loads and stores at a dummy block inside the
kernel frame with a zero advance stride, so inactive instances touch no memory
beyond the remainder and the hot loop stays branch-free.

\paragraph{\texttt{arm64}.}
The \texttt{arm64} path uses NEON plus the Armv8.2 SHA-3 extension
(\texttt{FEAT\_SHA3}), available on Apple Silicon and on Neoverse V1/V2 and N2.
The permutation is expressed with the extension's fused instructions:
\texttt{EOR3} (three-input XOR) for the $\theta$ column parities, \texttt{RAX1}
(rotate-by-one and XOR) for the $\theta$ $D$-terms, \texttt{XAR} (XOR then
rotate) to fuse the $\theta$ accumulation with the $\rho$ rotation, and
\texttt{BCAX} (bit-clear and XOR) for $\chi$. The single-duplex core places one
instance per 64-bit \texttt{D} lane across the 25 vector registers; the
eight-way core packs two instances per 128-bit register and processes the batch
of eight as four such pairs, using \texttt{ZIP1} and duplicating moves to
interleave and extract the two instances of each lane in the fused chunk kernel.
The same pair machinery backs a $2$-wide kernel that drains a leaf-chunk
remainder two chunks at a time---one instance per $64$-bit half of a $128$-bit
register---at roughly twice the single-duplex per-chunk rate. Both batched
kernels store the permutation state back into the batch after the closing
block---the $\times 8$ kernel for its first pair, the $2$-wide kernel for both
instances---which is the writeback that permits the root duplex to occupy
lane~$0$.
